%%%%%%%%%%  scan idx 296  |  folio 283  %%%%%%%%%%
\textbf{2.7. Loi contravariante et structure d'anneau sur \(K^{\bullet}\)} Nous allons expliciter quelques cas particuliers importants des accouplements (2.6) (leur étude directe est d'ailleurs immédiate).
\begin{enumerate}[label=\alph*)]
\item Soit \(f:(X',A')\longrightarrow (X,A)\) un morphisme de topos annelés. Le foncteur \(Lf^*:D^-({}_A X)\longrightarrow D^-({}_{A'} X')\) envoie (I 4.19.1) \(D({}_A X)_{\parf}\) dans \(D({}_{A'} X')_{\parf}\) (resp. \(\Loclib({}_A X)\) dans \(\Loclib({}_{A'} X')\)), et définit un homomorphisme
\begin{equation}\tag{2.7.1}
 f^*:K^{\bullet}({}_A X)\longrightarrow K^{\bullet}({}_{A'} X')
\end{equation}
(resp.
\begin{equation}\tag{2.7.2}
 f^*:K^{\bullet}({}_A X)_{\naif}\longrightarrow K^{\bullet}({}_{A'} X')_{\naif}
\end{equation}
), qu'on peut encore interpréter comme le produit par \(\cl(A')\) dans (2.6.5) (resp. (2.6.2)). On a un diagramme commutatif (cas particulier de (2.6.6))
\[
\begin{tikzcd}
K^{\bullet}({}_A X)_{\naif}\arrow[r,"f^*"]\arrow[d]
&K^{\bullet}({}_{A'} X')_{\naif}\arrow[d]\\
K^{\bullet}({}_A X)\arrow[r,"f^*"]
&K^{\bullet}({}_{A'} X'),
\end{tikzcd}
\]
où les flèches verticales de gauche et de droite sont respectivement (2.6.3.2) et (2.6.3.3).

Soit \(g:(X,A)\to(X'',A'')\) un morphisme de topos annelés, la transitivité des foncteurs images inverses implique que les homomorphismes du type (2.7.1) (resp. 2.7.2) vérifient la relation
\begin{equation}\tag{2.7.3}
 (gf)^*=f^*g^*.
\end{equation}
Ainsi \(K^{\bullet}({}_A X)\) (resp. \(K^{\bullet}({}_A X)_{\naif}\)) \emph{est un foncteur contravariant de} \((X,A)\).

\item Soit \((X,\mathcal O_X)\) un topos commutativement annelé. Grâce à l'associativité et à la commutativité du produit tensoriel, \emph{l'homomorphisme} (2.6.5) (resp. (2.6.2)) (où l'on fait \(A=\mathcal O_X\), \((X',A')=(X,A)\), \(f=\Id\)) \emph{munit} \(K^{\bullet}(X)\) (resp. \(K^{\bullet}(X)_{\naif}\)) \emph{d'une structure d'anneau commutatif avec unité} 
